% ==========================================================================
%  Paragraph: Runtime Scalability Analysis
% ==========================================================================

\paragraph{Runtime Scalability}
Figure~\ref{fig:elapsed_by_k} and Table~\ref{tab:elapsed_by_k}
present the mean elapsed time of the three algorithms as a function
of the goal count~$k$, averaged over all 2{,}500 problem instances.

Iterative-$k$A* is the fastest algorithm across the entire range of
$k$-values, with near-constant runtime between 1.6--2.5\,s
(slope $= 0.009$\,s per $k$-unit; growth factor $1.3{\times}$ from
$k{=}10$ to $k{=}100$). This confirms that sharing and reusing the
search effort across successive goal queries effectively amortizes
the per-goal cost.

In contrast, Aggregative-$k$A* exhibits strictly linear growth
(slope $= 0.082$\,s per $k$-unit), rising from 2.4\,s at $k{=}10$ to
9.3\,s at $k{=}100$ --- a $3.9{\times}$ increase. Since this approach
solves all goals in a single unified search, every additional goal
expands the open list and increases the overall exploration cost
proportionally.

$k$Dijkstra serves as an uninformed baseline with essentially
constant runtime (${\approx}6.4$\,s, slope $\approx 0$), as it
explores the entire reachable graph regardless of the number of goals.

A notable crossover occurs at $k \approx 65$
(Figure~\ref{fig:elapsed_by_k}): below this threshold,
Aggregative-$k$A* outperforms $k$Dijkstra by leveraging heuristic
guidance on a relatively small goal set; above it, the overhead of
managing a large aggregated open list exceeds the benefit of heuristic
pruning, and the uninformed exhaustive search of $k$Dijkstra becomes
preferable. At $k{=}100$, Iterative-$k$A* achieves a $3.9{\times}$
speedup over Aggregative-$k$A* and a $2.7{\times}$ speedup over
$k$Dijkstra, demonstrating that the iterative reuse strategy scales
gracefully where both alternatives degrade or stagnate.
